\chapter{扩展研究与综合分析}\label{chap:extension}
\section{扩展问题}
在基础方法之后，可以结合新的研究问题组织扩展章节。章节标题与内容应由论文实际贡献决定，模板提供的七章结构可以按需增减。

\section{误差来源与适用条件}
讨论误差来源时，应区分直接观察到的现象、作者确认的场景事实和仍待验证的原因。对实验设置不同的结果，应明确评价口径差异，并避免不具可比性的严格排名。

\subsection{长期运行条件}
长期运行实验可以从累计误差、观测覆盖和计算存储需求等方面讨论。论文应依据实际数据决定哪些分析具有意义，不为填满页面而加入缺少依据的结论。

\section{本章小结}
本章演示扩展分析的结构。若某章不适用于实际研究，可在主文件中删除对应的\verb|\input|，并同步更新章节安排文字。
